% GENERATED_BY: run_oc_core_monograph_translation_rework_dispatch_v1.py
% AI_ASSISTED_TRANSLATION_DRAFT: TRUE
% THEOREM_REVIEW_STATUS: APPROVED
% THEOREM_REVIEW_MODE: FORMAL_AI_GATE__CODEX_CLI_CHATGPT
% THEOREM_REVIEW_REVIEWER_ID: CODEX_CLI_CHATGPT__AUTO_DISPATCH_20260406
% THEOREM_REVIEW_AUTHORITY_CLASS: FINAL_ADVERSARIAL_EXTERNAL_LLM_REVIEWER
% THEOREM_REVIEWED_AT: 2026-04-14T12:38:35Z
% SOURCE_LANGUAGE: EN
% TARGET_LANGUAGE: RU
% SOURCE_EN_REF: content/frontmatter_oc_core_1_3_master.tex
\begin{titlepage}
\thispagestyle{empty}
\begin{center}
{\Large Онтология континуумов}\\[0.6cm]
{\huge \textbf{Core 1.3}}\\[0.3cm]
{\Large Мастер-монография с приоритетом первоисточника}\\[1.2cm]

{\Large Alexander Yashin}\\[0.2cm]
{\normalsize Независимый исследователь, Лейпциг/Галле, Германия}\\[0.2cm]
{\normalsize \href{https://orcid.org/0009-0008-6166-0914}{ORCID 0009-0008-6166-0914}}\\[1.0cm]

{\normalsize Версия 1.3}\\[0.2cm]
{\normalsize 13 апреля 2026 г.}\\[1.2cm]

\begin{minipage}{0.88\textwidth}
\small
\textbf{Публикационная роль.} Этот том является канонической мастер-монографией для OC Core 1.3.
Он восстанавливает полную научную оболочку монографического масштаба, сохраняет унаследованный формальный корпус
Core 1.2 как субстрат в форме, родной для первоисточника, и добавляет явные слои аудита источника, судьбы теорем,
хронологии, дорожной карты теорем, разобранных примеров и границ публикации, требуемые для релиза 1.3.

\medskip

\textbf{Обещание читателю.} Цель этого документа не в том, чтобы сжать каркас в
краткий пакет. Цель состоит в том, чтобы дать читателю один проверяемый научный том, в котором
определения, формулировки теорем, доказательства, рисунки, ограничения и происхождение материала можно читать в
их естественном порядке, не вынуждая читателя восстанавливать теорию по артефактам дашбордов
или скрытым сопутствующим реестрам. Поэтому том включает явный путеводитель для читателя,
дорожную карту теорем, разобранные примеры и приложение для навигации рецензента, чтобы
враждебный рецензент и серьёзный неспециалист могли войти в один и тот же научный объект по разным допустимым
путям, не изменяя границу доказательства.
\end{minipage}

\vfill

% DEDICATION_REQUIRED_DO_NOT_REMOVE
% OC_CORE_PUBLICATION_DEDICATION
{\small\itshape Посвящается моей дорогой жене Марии, без которой эта работа была бы невозможной.}\\[0.8cm]

{\small Открытый маршрут основы: github.com/alexanderyashin/ontology-of-continua-core-main}\\
{\small Архивное свидетельство предшественника: Zenodo DOI 10.5281/zenodo.17903912}
\end{center}
\end{titlepage}

\begin{abstract}
Core 1.3 является повторно валидированной мастер-монографией «Онтологии континуумов» с приоритетом первоисточника.
Она исходит из полной формальной оболочки и дерева глав, которые сделали Core 1.2 первым
аксиоматически замкнутым релизом OC, а затем подчиняет этот унаследованный корпус более строгой публикационной
дисциплине: явному засвидетельствованию источника, разделению по судьбе теорем, примечаниям об исправлениях с учётом
хронологии, границам незаявляемого и академической упаковке, готовой к внешней публикации.

Этот результат намеренно больше, чем извлечённая из него короткая журнальная статья.
Эта монография является дидактическим и формальным мастер-источником: она несёт структурную теорию,
аксиоматический и теоремный каркас, иерархию континуумов и пространств вложения,
межуровневые и предметные модули, программы фальсифицируемости и явные материалы 1.3,
которые объясняют, что сохраняется без изменений, что пересматривается, что ограничивается, а что откладывается.

Поэтому научная роль этого тома двояка. Во-первых, он сохраняет читаемость монографического уровня,
которой требует серьёзный фундаментальный каркас. Во-вторых, он предоставляет канонический источник, из которого
можно выводить более узкие журнально-ядерные и канал-специфические артефакты без
молчаливого расширения или сужения защищаемого утверждения. Тем самым монография является как самостоятельной
научной работой, так и авторитетным субстратом для внешней публикации, рецензирования, перевода
и последующей проекции в продуктовые, бизнес- и публичные каналы.

Релиз 1.3 также делает навигацию по доказательствам явной. Вместо предположения, что каждый читатель
сам обнаружит один и тот же путь через большую формальную оболочку, том теперь включает отдельную
дорожную карту теорем, разобранные примеры семантики коллапса / остатка / возрождения и доктрины уровня K, а также
матрицу навигации для рецензента, которая сводит возражения, цели доказательств и свидетельства источника
в один проверяемый маршрут.
\end{abstract}

\clearpage
